% Part: computability
% Chapter: recursive-functions
% Section: primitive-recursion

\documentclass[../../../include/open-logic-section]{subfiles}

\begin{document}

\olfileid{cmp}{rec}{pre}
\olsection{بنسټيز بازګښت}

د طبيعي عددونو يو خاصيت دا دے چې هر طبيعي عدد پر~$0$ د راتلونکي
عمل~$+1$ په متناهي څو ځله کارولو سره تر لاسه کېدے شي---هر طبيعي عدد
يا~$0$ دے، يا د \dots{} د~$0$ راتلونکی دے۔ د دې حقيقت په کارولو سره
د يوې تابع~$h\colon\Nat \to \Nat$ د ټاکلو يوه لاره داسې ده:
(الف)~وټاکئ چې د آرګومېنټ~$0$ دپاره د~$h$ ارزښت څۀ دے، او
(ب)~دا هم وټاکئ چې د~$h(x)$ ارزښت له ورکړې وروسته د~$h(x+1)$ ارزښت
څنګه محاسبه شي۔ ځکه (الف) په نېغه راته وايي چې~$h(0)$ څۀ دے، نو
$h$ د~$0$ دپاره تعريف شوه۔ اوس د~$x=0$ دپاره د (ب) لارښوونه په
کارولو سره، له~$h(0)$ څخه $h(1) = h(0+1)$ محاسبه کولے شو۔ هماغه
لارښوونه د~$x=1$ دپاره په کارولو سره، له~$h(1)$ څخه
$h(2) = h(1+1)$ محاسبه کوو، او همداسې مخکښې ځو۔ د هر طبيعي
عدد~$x$ دپاره به بالاخره هغه پړاو ته ورسېږو چې~$h(x+1)$ له~$h(x)$
څخه تعريفوو؛ نو~$h(x)$ د ټولو~$x \in \Nat$ دپاره تعريف شوې ده۔
\textbf{د ژباړې د نسخې سمون (OLCMP-002):} سرچينې د دې جملې د
بازګښت لوری سرچپه ليکلی و او اوسنی ارزښت ئې له راتلونکي ارزښت څخه
تعريفاوه؛ دلته د مخکښې ورکړې لارښوونې له مخې راتلونکی ارزښت له اوسني
ارزښت څخه تعريف شوے دے۔

د بېلګې په توګه، فرض کړئ~$h\colon \Nat \to \Nat$ په لاندې دوو
معادلو ټاکو:
\begin{align*}
h(0) & =  1\\
h(x+1) & =  2 \cdot h(x)
\end{align*}
که د ضرب طريقه له مخکښې را معلومه وي، نو دا معادلې د پورته (الف)
او~(ب) دپاره اړين معلومات راکوي۔ د دويمې معادلې پرله‌پسې په کارولو
سره لاس ته راوړو چې
\begin{align*}
  h(1) & = 2\cdot h(0) = 2,\\
  h(2) & = 2\cdot h(1) = 2\cdot 2,\\
  h(3) & = 2 \cdot h(2) = 2\cdot 2 \cdot 2,\\
  & \vdots
\end{align*}
وينو چې هغه تابع~$h$ چې ټاکلې مو ده، $h(x) = 2^x$ ده۔

د طبيعي عددونو ځانګړے خاصيت ضمانت کوي چې يوازې يوه تابع~$h$ شته
چې دا دواړه معيارونه پوره کوي۔ د دې غوندې معادلو جوړې ته د
تابع~$h$ \emph{په بنسټيز بازګښت تعريف} وايي۔ دې ته ځکه
«بازګشتي» وايي چې~$h$ په بازګشتي ډول تعريفوو، يعنې تعريف، او په
ځانګړي ډول دويمه معادله، پخپله~$h$
په ښي اړخ کښې لري۔ دې ته ځکه «بنسټيز» وايي چې د~$h(x+1)$ په تعريف
کښې يوازې د~$h(x)$ ارزښت، يعنې سملاسي مخکينی ارزښت، کاروو۔ دا
پر~$\Nat$ د يوې تابع د بازګشتي تعريف تر ټولو ساده لاره ده۔

د جمع او ضرب غوندې لا بنسټيزې تابعې هم په بنسټيز بازګښت تعريفولے
شو۔ خو په دې حالتونو کښې اړوندې تابعې~$2$-ځاييزې دي۔ د آرګومېنټونو
يو ځاے ثابت نيسو او بل د بازګښت دپاره کاروو۔ د بېلګې په توګه، د
$\Add(x, y)$ د تعريف دپاره~$x$ ثابت نيولے شو، ارزښت لومړے د~$y=0$
دپاره تعريفولے شو، او بيا د~$y+1$ دپاره ئې د~$y$ له مخې تعريفولے
شو۔ ځکه چې~$x$ ثابت دے، د تعريفي معادلو په کيڼ او ښي اړخ دواړو
کښې به راځي۔
\begin{align*}
\Add(x,0) & =  x\\
\Add(x,y+1) & =  \Add(x,y)+1
\end{align*}
دا معادلې د ټولو~$x$ \emph{او}~$y$ دپاره د~$\Add$ ارزښت ټاکي۔ د
بېلګې په توګه، د~$\Add(2,3)$ د موندلو دپاره تعريفي معادلې د~$x = 2$
دپاره کاروو: لومړۍ معادله کاروو چې~$\Add(2,0) = 2$ ومومو، بيا
دويمه معادله پرله‌پسې کاروو چې~$\Add(2,1) = 2 + 1 = 3$،
$\Add(2, 2) = 3 + 1 = 4$، او~$\Add(2, 3) = 4 + 1 = 5$ ومومو۔

د~$\Add$ په تعريف کښې مو د دويمې معادلې په ښي اړخ کښې~$+$ وکاراوه،
خو يوازې د~$1$ د زياتولو دپاره۔ په بله وينا، د راتلونکي
تابع~$\Succ(z) = z+1$ مو وکاروله او پر مخکيني ارزښت~$\Add(x,y)$ مو
عملي کړه، څو~$\Add(x,y+1)$ تعريف کړو۔ نو بازګشتي تعريف داسې ګڼلے شو
چې د يوې تابع له مخې ورکړل شوے دے او موږ ئې پر مخکيني ارزښت عملي
کوو۔ خو که تابع يوازې پر مخکيني ارزښت نۀ، بلکې پر~$x$ او~$y$ هم
تابع وګرځوو، زيان نۀ لري---او کله ناکله اړينه وي۔ دا تعريف وګورئ:
\begin{align*}
  \Mult(x,0) & =  0 \\
  \Mult(x,y+1) & =  \Add(\Mult(x,y),x)
\end{align*}
دا د تابع~$\Mult$ يو بنسټيز بازګشتي تعريف دے چې تابع~$\Add$ هم پر
مخکيني ارزښت~$\Mult(x,y)$ او هم پر لومړي آرګومېنټ~$x$ عملي کوي۔
دا هم تابع~$\Mult(x,y)$ د ټولو آرګومېنټونو~$x$ او~$y$ دپاره
تعريفوي۔ د بېلګې په توګه، $\Mult(2,3)$ د~$\Mult(2,0)$،
$\Mult(2,1)$، $\Mult(2,2)$ او~$\Mult(2,3)$ په پرله‌پسې محاسبې
سره ټاکل کېږي:
\begin{align*}
  \Mult(2,0) & = 0\\
  \Mult(2,1) & = \Mult(2,0+1) =
  \Add(\Mult(2,0), 2) = \Add(0, 2) = 2\\
  \Mult(2,2) & = \Mult(2,1+1) =
  \Add(\Mult(2,1), 2) = \Add(2, 2) = 4\\
  \Mult(2,3) & = \Mult(2,2+1) =
  \Add(\Mult(2,2), 2) = \Add(4, 2) = 6
\end{align*}

نو عمومي نمونه داسې ده: د يوې تابع~$h(x_0, \dots, x_{k-1}, y)$
د بنسټيز بازګشتي تعريف دپاره دوه معادلې ورکوو۔ لومړۍ د~$h(x_0,
\dots, x_{k-1}, 0)$ ارزښت~$h$ ته له اشارې پرته تعريفوي۔ دويمه د
$h(x_0, \dots, x_{k-1}, y+1)$ ارزښت د~$h(x_0, \dots,
x_{k-1}, y)$، نورو آرګومېنټونو~$x_0$، \dots،~$x_{k-1}$، او~$y$
له مخې تعريفوي۔ په دې دويمه معادله کښې يوازې د~$h$ سملاسي مخکينی
ارزښت کارېدے شي۔ که د دې دوو معادلو د ښي اړخونو ورکړي عملونه هم
تابعې~$f$ او~$g$ وګڼو، نو په بنسټيز بازګښت د نوې تابع~$h$ د تعريف
عمومي نمونه داسې ده:
\begin{align*}
  h(x_0, \dots, x_{k-1}, 0) & = f(x_0, \dots, x_{k-1})\\
  h(x_0, \dots, x_{k-1}, y+1) & =
  g(x_0, \dots, x_{k-1}, y, h(x_0, \dots, x_{k-1}, y))
\end{align*}
د~$\Add$ په حالت کښې، $k=1$ او~$f(x_0) = x_0$ (عينيت تابع) دي،
او~$g(x_0, y, z) = z + 1$ (هغه~$3$-ځاييزه تابع چې د خپل درېيم
آرګومېنټ راتلونکی ورکوي) ده:
\begin{align*}
  \Add(x_0, 0) & = f(x_0) = x_0\\
  \Add(x_0, y+1) & = g(x_0, y, \Add(x_0, y)) =
  \Succ(\Add(x_0, y))
\end{align*}
د~$\Mult$ په حالت کښې، $f(x_0) = 0$ (هغه ثابت تابع چې تل~$0$
ورکوي) او~$g(x_0, y, z) = \Add(z,x_0)$ (هغه~$3$-ځاييزه تابع چې د
خپل وروستي او لومړي آرګومېنټ مجموعه ورکوي) دي:
\begin{align*}
  \Mult(x_0,0) & =  f(x_0) = 0 \\
  \Mult(x_0,y+1) & =  g(x_0, y, \Mult(x_0,y)) =
  \Add(\Mult(x_0,y), x_0)
\end{align*}

\end{document}
